% Introduction — ML contribution drop-in fragment
% Author: ML (agent-mozurrmd). Source: workspace/proposals/unified_ML_spec_v2.md §4.
% Drop into sections/01_introduction.tex after the problem-statement paragraph.
% Biosci: feel free to trim or reorder; citations follow the lab bibfile conventions.

Over the past decade, single-cell batch integration has evolved along four broadly separable lines. The first is parametric probabilistic modeling of expression as a batch-conditioned generative process — scVI \citep{lopez2018scvi}, scANVI \citep{xu2021scanvi}, scPoli \citep{hrovatin2023scpoli}, MultiVI, totalVI — where batch is a covariate and the learned latent is argued, but not demonstrated, to be biology-only. The second is anchor- and contrastive-based alignment — Seurat v3--v5 \citep{stuart2019seurat, hao2021seurat}, fastMNN \citep{haghverdi2018mnn}, Harmony \citep{korsunsky2019harmony}, Scanorama \citep{hie2019scanorama}, scDML \citep{yu2023scdml}, scDREAMER \citep{shaham2023scdreamer}, Concerto \citep{yang2022concerto}, INSCT \citep{simon2021insct} — which rely on mutual-nearest-neighbor or cluster-centroid primitives that, by construction, rigidify the embedding along dense cross-batch correspondences. The third is semi-supervised and label-informed integration — scANVI, STACAS \citep{andreatta2024stacas}, scPoli — where prior labels bias the objective. The fourth is post-hoc signal recovery — CellANOVA \citep{li2024cellanova} — which retrofits lost variation after integration rather than preventing its loss.

Three developments from 2022--2026 reframe the field. SCIDRL \citep{wang2022scidrl}, scCRAFT \citep{li2025sccraft}, and RBET \citep{wang2025rbet} each target rare-subpopulation preservation explicitly. scIB-E \citep{xu2025scibe} shows that within-cell-type biological variation is invisible to the scIB scoring system \citep{luecken2022scib}. A February-2026 review in \textit{Nature Computational Science} summarizes that, despite the algorithmic diversity above, overcorrection remains unresolved. The most recent semi-supervised benchmark \citep{plos2025semisupbench} shows that even label-informed methods collapse under mild label noise.

A single epistemic invariant underlies every result in this list: evaluation is label-bound. scDML, scDREAMER, scCRAFT, and SCIDRL validate \emph{preservation of annotated} rare cells; CellANOVA recovers \emph{known} biological condition differences; scIB, scIB-E, and RBET score clustering against \emph{external} labels. As the measure-theoretic argument in Methods (Theorem~\ref{thm:label-blindness}) makes rigorous, \emph{no label-bound evaluation functional is sensitive to local perturbations that destroy an unannotated rare component.} The algorithmic consequence is a systematic bias whose size is bounded below by the rarity of the unannotated population --- exactly the regime our framework targets.

\textbf{REAL} sits orthogonally. It adopts the transport-theoretic lens initiated by Waddington-OT \citep{schiebinger2019wot}, Unbalanced OT \citep{chizat2018uot}, SCOT \citep{demetci2020scot}, and MIOFlow \citep{bunne2022monge}, but for the first time uses it as an \emph{evaluation} signal rather than an integration algorithm. It inherits the density-preservation stance of manifold learners \citep{moon2019phate, moor2020topoae} and of the OOD / open-set literature \citep{bendale2016openmax, lee2018mahalanobis, liu2020energy, nalisnick2019typicality, ren2019likelihood}, bringing per-neighborhood typicality scoring to the integration setting. It addresses the long-tail / imbalanced representation problem \citep{menon2021logitadj, cui2019classbalanced} by treating rare witnesses as a sparse structured signal rather than per-sample class imbalance. And it replies to domain-adaptation overcorrection \citep{bendavid2010domain, ganin2016dann} by explicitly exempting witnessed low-density regions from domain-invariance penalties.

The result is a framework that, unlike any prior method, (i)~detects unannotated rare-subpopulation loss without labels, (ii)~protects against it during training by minimizing the same inequality the detector thresholds, and (iii)~certifies, via a conservation-of-local-mass identifiability bound, under what conditions detection is statistically possible.
